\reading{LTR-16}{Covered Interest Parity Lost: Understanding the Cross-Currency Basis}{QUICK WIN}{2}
\markboth{LTR-16\ \ Covered Interest Parity Lost}{}

\begin{lrkeybox}[Learning objectives — official 2026 spine wording]
\begin{enumerate}[label=\textbf{\alph*.}]
\item Differentiate between the mechanics of foreign exchange (FX) swaps and cross-currency swaps.
\item Identify key factors that affect the cross-currency swap basis.
\item Assess the causes of covered interest rate parity violations after the financial crisis of 2008.
\end{enumerate}
{\footnotesize\color{FRMGrey}Source: Claudio Borio, Robert McCauley, Patrick McGuire and Vladyslav Sushko, ``Covered Interest Parity Lost: Understanding the Cross-Currency Basis,'' BIS Quarterly Review, Q3 2016.}
\end{lrkeybox}

\begin{lrtrapbox}[Where this reading's first page went]
The Lecture layer's opening page for this reading sits under the heading ``\texttt{LO 72.g: Calculate the financing advantage of a bond trading special}'' --- an orphan heading inherited from LTR-13. Its content is the Covered Interest Parity material below. Reunited here; see the corrections record and the note at LTR-13 g.
\end{lrtrapbox}

% =====================================================================
\lo{LTR-16 a}{Differentiate between the mechanics of foreign exchange (FX) swaps and cross-currency swaps.}

\begin{lrdefbox}[Covered interest parity]
A no-arbitrage condition according to which interest rates on two otherwise identical assets in two different currencies should be equal \textbf{once the foreign currency risk is hedged}.
\end{lrdefbox}

\begin{lrfmlbox}[CIP, in two equivalent forms]
\[ \frac{F}{S} \;=\; \frac{1+r}{1+r^{*}} \qquad\qquad\qquad F - S \;=\; S\left(\frac{1+r}{1+r^{*}} - 1\right) \]
\vk{$S$ = spot exchange rate in US dollars per unit of foreign currency (USD/FC); $F$ = forward exchange rate in USD/FC; $r$ = the US dollar interest rate; $r^{*}$ = the foreign currency interest rate. The second form says the \textbf{forward premium} $(F-S)$, in points, should equal the spot rate times the interest rate differential.}
\end{lrfmlbox}

\begin{lrkeybox}[The intuition: two ways to invest in US dollars]
\begin{enumerate}
\item \textbf{Directly}, by investing at the USD interest rate $r$.
\item \textbf{Indirectly}, by exchanging USD for FC at the spot rate $S$, investing the FC at $r^{*}$, and locking in the return by agreeing to convert the FC proceeds back into USD at maturity at the forward rate $F$.
\end{enumerate}
These two trades should have the same return, or a risk-free arbitrage opportunity is available.
\end{lrkeybox}

\subsection{FX swap versus cross-currency swap}

\begin{center}
\small
\begin{tabularx}{\textwidth}{@{}p{2.5cm} X X@{}}
\toprule
\rowcolor{FRMSoft}\textbf{Feature} & \textbf{FX swap} & \textbf{Cross-currency swap} \\
\midrule
\textbf{Mechanics} &
One party \textbf{borrows one currency from}, and \textbf{simultaneously lends another currency to}, a second party. &
The two parties \textbf{also simultaneously borrow and lend} an equivalent amount of funds in two different currencies. \\
\addlinespace
\textbf{Tenor} & Short-term & \textbf{Longer-term}, typically \textbf{above one year} \\
\addlinespace
\textbf{Exchange at maturity} &
Borrowed amounts are exchanged at the \textbf{spot} rate at inception and repaid at the \textbf{pre-agreed forward} rate at maturity. &
Borrowed amounts are exchanged back at the \textbf{initial spot rate}. \\
\addlinespace
\textbf{Interim payments} & None. & Counterparties \textbf{periodically exchange interest payments} during the life of the swap. \\
\addlinespace
\textbf{Purpose} & A way of funding an asset denominated in a foreign currency by paying interest in the \textbf{domestic} currency. & Longer-horizon funding and hedging. \\
\bottomrule
\end{tabularx}
\end{center}
\figcap{Two discriminators carry the objective. \textbf{Tenor}: an FX swap is short-dated, a cross-currency swap runs beyond a year. \textbf{Re-exchange rate}: the FX swap unwinds at the \emph{forward} rate agreed at inception, the cross-currency swap unwinds at the \emph{initial spot} rate --- and that is possible only because the cross-currency swap pays interest along the way.}

\begin{lrtrapbox}[Forward rate versus initial spot rate at maturity]
This is the cleanest single-word discriminator in the reading. \textbf{FX swap $\rightarrow$ pre-agreed forward rate. Cross-currency swap $\rightarrow$ initial spot rate.} Swapping the two produces a statement that reads perfectly and is wrong, and it is the obvious \S5B.2 sibling build.
\end{lrtrapbox}

\subsection{The cross-currency basis}

Since 2008 the USD has tended to \textbf{command a premium} relative to the FC in FX swaps: the party \emph{lending} USD can sell the FC forward at a price $F$ that is \textbf{higher} than indicated by the interest rate differential. Typically $(F - S)$ has been \textbf{``too wide''} and CIP has not held, such that:

\begin{lrfmlbox}[The CIP violation, and the basis that restores it]
\[ F - S \;>\; S\left(\frac{1+r}{1+r^{*}} - 1\right) \qquad\text{(CIP fails, $F-S$ too wide)} \]
\[ F - S \;=\; S\left(\frac{1 + r + b}{1 + r^{*}}\right) - S \qquad\text{(the basis $b$ restores equality)} \]
\vk{$b$ = the \term{cross-currency swap basis}: the amount by which the \textbf{USD interest rate} must be adjusted so that CIP holds, assuming $F-S$ is too wide. Note where $b$ sits --- it is added to $r$ in the \emph{numerator}, alongside the dollar rate.}
\end{lrfmlbox}

\begin{lrkeybox}[Who pays the basis]
When CIP does not hold, the party \textbf{borrowing the USD pays the basis \emph{in addition to} the USD interest rate}, while the party \textbf{borrowing the FC pays \emph{less than} the FC interest rate} by the amount of the basis.

If CIP holds throughout the life of the contract, then \textbf{the basis is zero} and the cross-currency basis swap is simply a straightforward floating-for-floating currency swap.
\end{lrkeybox}

\begin{lrtrapbox}[The basis is a charge on the dollar borrower]
Two polarity facts move together: the dollar borrower \textbf{pays more}, the foreign-currency borrower \textbf{pays less}. Reversing either half is a single-word corruption of a true statement. And when $b = 0$ the instrument reduces to an ordinary floating-for-floating currency swap --- an option describing a zero basis as making the swap impossible or meaningless has the wrong conclusion.
\end{lrtrapbox}

% =====================================================================
\lo{LTR-16 b}{Identify key factors that affect the cross-currency swap basis.}

Two key factors influence the cross-currency swap basis and can lead to deviations from CIP.

\begin{center}
\small
\begin{tabularx}{\textwidth}{@{}p{3.6cm} X@{}}
\toprule
\rowcolor{FRMSoft}\textbf{Factor} & \textbf{Mechanism} \\
\midrule
\textbf{1. Decline in market liquidity} &
During financial crises, market liquidity in the underlying \textbf{forward and spot FX contracts} can decrease significantly. Reduced liquidity leads to \textbf{wider bid-ask spreads}, increasing the transaction cost of executing FX swaps and making arbitrage less feasible. When crises hit, FX markets dry up, making trades expensive and eliminating the arbitrage that keeps the basis near zero. \\
\addlinespace
\textbf{2. Risk premia} &
Increases in \textbf{counterparty credit risk} --- measured by \textbf{LIBOR--OIS spreads} --- in FX markets, and \textbf{sovereign credit risk} --- measured by \textbf{sovereign CDS spreads} --- in the government bonds used for CIP arbitrage, result in risk premia. \textbf{Even small increases} in these risk premia can hinder arbitrage execution, particularly when there is substantial hedging demand for one of the currencies involved. \\
\bottomrule
\end{tabularx}
\end{center}

\begin{lrtrapbox}[Two risks, two named measures]
\textbf{Counterparty credit risk $\rightarrow$ LIBOR--OIS spread. Sovereign credit risk $\rightarrow$ sovereign CDS spread.} The measures are named specifically and are swappable in a distractor. Note also the qualifier the reading insists on: \textbf{even small} increases are enough, because the arbitrage margin is thin to begin with.
\end{lrtrapbox}

% =====================================================================
\lo{LTR-16 c}{Assess the causes of covered interest rate parity violations after the financial crisis of 2008.}

The reading gives two halves: why the basis \emph{opens up}, and why it \emph{does not close}.

\subsection{Why the basis opens up: demand for currency hedges}

Two factors affect the cross-currency swap basis: the \textbf{risk premium} for the underlying investment over the duration of the swap, and \textbf{hedging demand}. Three sources of hedging demand are \textbf{rather insensitive to the size of the basis}, and hence exert \textbf{sustained pressure on it even when it is non-zero}:

\begin{center}
\small
\begin{tabularx}{\textwidth}{@{}p{4.8cm} X@{}}
\toprule
\rowcolor{FRMSoft}\textbf{Source of hedging demand} & \textbf{What they do} \\
\midrule
\textbf{Banks managing currency mismatches} & Use FX swaps to hedge currency mismatches on their balance sheets. \\
\addlinespace
\textbf{Institutional investors hedging portfolios} & Insurance companies and banks use FX swaps to hedge foreign currency risk in their investment portfolios. \\
\addlinespace
\textbf{US firms issuing foreign-currency debt} & US nonfinancial companies issue foreign-currency debt and then \textbf{swap the proceeds back into USD} using FX swaps. \\
\bottomrule
\end{tabularx}
\end{center}

\begin{lrtrapbox}[Price-insensitive demand is the whole mechanism]
The examinable property of these three is that they are \textbf{insensitive to the size of the basis}. Ordinary demand recedes when the price moves against it and the market clears; demand that does not recede keeps pressing on the basis indefinitely. An option describing hedging demand as basis-sensitive, or as receding once the basis widens, removes the reason the basis persists.
\end{lrtrapbox}

\subsection{Limits to arbitrage: why the basis does not close}

\begin{itemize}
\item Arbitrage became \textbf{both costly and risky after the crisis} --- including in emerging markets with capital restrictions.
\item As a result of \textbf{tighter management of risks} and related \textbf{balance sheet constraints}, arbitrage now incurs a \textbf{cost per unit of balance sheet}. Balance sheet deleveraging compounds this.
\item That cost is \textbf{passed on to the pricing of FX swaps}, introducing a premium.
\item \textbf{Changes in regulation} have reinforced market pressures for tighter management of balance sheet risks.
\end{itemize}

\begin{lrtrapbox}[The cost is per unit of balance sheet, not per trade]
This is the phrase the reading turns on. A CIP arbitrage is riskless in payoff terms but is \emph{not} free, because it consumes balance sheet capacity that is now scarce and charged for. That is why a violation can persist without anyone being irrational: the arbitrage exists, and executing it costs more than it pays.
\end{lrtrapbox}

\begin{lrtrapbox}[Consolidated trap summary — LTR-16]
\begin{itemize}
\item \textbf{Sibling.} FX swap unwinds at the \textbf{pre-agreed forward rate}; cross-currency swap unwinds at the \textbf{initial spot rate} and pays interest in between.
\item \textbf{Scope.} Cross-currency swaps are \textbf{longer-term, typically above one year}.
\item \textbf{Polarity.} The \textbf{USD borrower pays the basis on top of} the USD rate; the \textbf{FC borrower pays less than} the FC rate by the same amount.
\item \textbf{Scope.} A \textbf{zero basis} means CIP holds and the instrument is an ordinary floating-for-floating currency swap.
\item \textbf{Definition.} The basis $b$ adjusts the \textbf{USD interest rate} in the numerator, not the spot or forward rate.
\item \textbf{Sibling.} Counterparty credit risk $\rightarrow$ \textbf{LIBOR--OIS}; sovereign credit risk $\rightarrow$ \textbf{sovereign CDS}.
\item \textbf{Scope.} The three hedging-demand sources are \textbf{insensitive to the size of the basis}, which is why the pressure is sustained.
\item \textbf{Scope.} Post-crisis arbitrage carries a \textbf{cost per unit of balance sheet}, passed into FX swap pricing.
\item \textbf{Polarity.} Since 2008 the USD commands a \textbf{premium}: $(F-S)$ has been \textbf{too wide}.
\end{itemize}
\end{lrtrapbox}
